\documentclass[french]{article}
\usepackage[utf8]{inputenc}
\usepackage[T1]{fontenc}
\usepackage{lmodern}
\usepackage{geometry}
\usepackage{graphicx}
\usepackage{babel}
\usepackage{amsmath}
\usepackage{amsthm}
\usepackage{amssymb}
\usepackage{hyperref}
\usepackage{stmaryrd}
\usepackage{enumitem}
\usepackage{tcolorbox}
\usepackage{dsfont}
\usepackage{multirow}
\usepackage{etoc}
\usepackage{titlesec}
\usepackage{esint}
\usepackage{cancel}
\usepackage{mathdots}

\setlist[itemize]{label=$\bullet$}


\renewcommand{\P}{\mathbb{P}}
\newcommand{\N}{\mathbb{N}}
\newcommand{\Z}{\mathbb{Z}}
\newcommand{\Q}{\mathbb{Q}}
\newcommand{\R}{\mathbb{R}}
\newcommand{\C}{\mathbb{C}}
\newcommand{\K}{\mathbb{K}}
\renewcommand{\L}{\mathbb{L}}
\newcommand{\U}{\mathbb{U}}
\newcommand{\st}{^*}
\newcommand{\tq}{\middle|}
\newcommand{\inv}{^{-1}}
\newcommand{\E}{\mathbb{E}}
\newcommand{\F}{\mathbb{F}}
\newcommand{\V}{\mathbb{V}}
\renewcommand{\ker}{\text{Ker}}
\newcommand{\im}{\text{Im}}
\newcommand{\card}{\text{Card}}
\newcommand{\Tr}{\text{Tr}}

\title{TD Euclidiens\\ Exercice 30}
\date{}
\begin{document}
\maketitle

\section{Énoncé}
Soit $E$ un espace euclidien et $u\in\mathcal{L}(E)$ tel que $\forall x\in E,\ \lVert u(x)\rVert\leqslant\lVert x\rVert$. Pour $x\in E$ et $n\in\N$, on pose : $$v_n(x)=\frac{1}{n+1}\sum_{k=0}^{n}u^k(x)$$
\begin{enumerate}
	\item Montrer, pour tout $x$ appartenant à $\ker(u-\text{Id}_E)$ ou à $\text{Im}(u-\text{Id}_E)$, que la suite $(v_n(x))_{n\in\N}$ converge, et préciser sa limite.
	\item En déduire $E=\ker(u-\text{Id}_E)\oplus \text{Im}(u-\text{Id}_E)$ et que pour tout $x\in E$, $(v_n(x))_{n\in\N}$ converge.
	\item $\star$ Montrer que cette somme directe est orthogonale.
\end{enumerate}
\section{Solution}
On peut construire $$v_n=\frac{1}{n+1}\sum_{k=0}^{n}u^k(x)$$ qui vaut $x$ sur $\text{Ker}(u-\text{id})$ et tend vers 0 sur $\text{Im}(u-\text{id})$. Cela permet de montrer que $$E=\text{Ker}(u-\text{id})\oplus\text{Im}(u-\text{id})$$ $v_n$ tend alors vers la projection sur $\text{Ker}(u-\text{id})$ parallèlement à $\text{Im}(u-\text{id})$. De plus, $v_n$ est de norme triple inférieure à $1$, donc la projection vers laquelle elle tend aussi. Cette projection est donc une projection orthogonale (\textit{cf} une caractérisation des projecteurs orthogonaux), donc la somme directe prouvée précédemment est orthogonale.

\end{document}
